\section{Auxiliary polynomials and local estimates}\label{sec:local}

Passing to an arithmetic progression changes the polynomial whose values are
forbidden differences. We first show that one such change can make the first
nonzero Taylor coefficient at each chosen local root a unit. We then organize
all subsequent changes into a coherent family \(h_l\), indexed by positive
integers, and show that this local property is preserved. The family and its
local estimates remain available for an arbitrary fixed intersective
polynomial; the unit condition will let us choose their structural parameters
using only the degree. The estimates give decay of complete exponential sums
at large primes, together with residue classes on which polynomial values
lift uniformly and support cycles of one fixed length.

Throughout this section, \(h\in\Z[x]\) is the fixed intersective polynomial of
degree \(k\ge2\), with positive leading coefficient \(a\).

\subsection{A family compatible with progressions}

We use Lucier's auxiliary-polynomial construction
\cite[Section~8]{Lucier2006}, in which a compatible root at each prime
determines the normalization by its multiplicity. We give the construction
and its inheritance rule explicitly, before proving the uniform local
estimates needed here.

Here \(\Z_p\) denotes the ring of compatible residues modulo the powers of
a prime \(p\), and \(v_p\) its valuation, with \(v_p(0)=\infty\).
Every nonzero element is a power of \(p\) times a unit, so \(\Z_p\) is an
integral domain. Intersectivity supplies a root of \(h\) modulo every power
of \(p\). These roots can be chosen compatibly: at each stage, retain a
residue class having root extensions to arbitrarily high powers, and then
choose one of its finitely many refinements with the same property. Fix the
resulting root \(\mathfrak z_p\in\Z_p\) for each prime \(p\), and write
\[
 h(\mathfrak z_p+X)=\sum_{j=m_p}^k b_{j,p}X^j,
 \qquad 1\le m_p\le k,
 \qquad b_{m_p,p}\ne0.
\]
Thus \(m_p\) is the multiplicity of the chosen root.

The coefficient \(b_{m_p,p}\) is a unit for all sufficiently large primes,
with a threshold depending only on \(h\). To see this, factor
\(h=c\prod_i g_i^{d_i}\), where \(c\) is a nonzero integer and the \(g_i\)
are pairwise nonassociate primitive irreducible polynomials in \(\Z[x]\). Cleared
B\'ezout identities for \((g_i,g_j)\), \(i\ne j\), and for \((g_i,g_i')\)
have nonzero integer right-hand sides. Exactly one factor \(g_i\) vanishes
at \(\mathfrak z_p\), and \(g_i'(\mathfrak z_p)\ne0\). Outside the finite
set of primes dividing these right-hand sides or \(c\), the derivative and
all other factor values at \(\mathfrak z_p\) are units. Hence \(m_p=d_i\)
and
\[
 b_{m_p,p}
 =c\,g_i'(\mathfrak z_p)^{d_i}
       \prod_{j\ne i}g_j(\mathfrak z_p)^{d_j}
\]
is a unit. This argument allows repeated factors of \(h\).

We will sometimes impose the following additional condition on the chosen
roots:
\begin{equation*}\tag{U}\label{local:unit-condition}
 b_{m_p,p}\in\Z_p^\times\qquad\text{for every prime }p.
\end{equation*}
The construction below works without this condition. Its role is to make the
local parameter choices depend only on \(k\). Every intersective polynomial
can be reduced to this case by one progression change.

\begin{lemma}[Normalization at the chosen roots]\label{local:normalization}
There are positive integers \(D_0,M\), an integer \(-D_0<r\le0\), and an
intersective polynomial
\[
 g(x)=\frac{h(r+D_0x)}{M}\in\Z[x]
\]
of degree \(k\) with positive leading coefficient such that, for every prime
\(p\), the element \(w_p=(\mathfrak z_p-r)/D_0\) belongs to \(\Z_p\) and
\[
 g(w_p+X)=\sum_{j=m_p}^k b^*_{j,p}X^j,
 \qquad b^*_{j,p}\in\Z_p,
 \qquad b^*_{m_p,p}\in\Z_p^\times.
\]
Thus \(g\), with these chosen roots, satisfies
Condition~\eqref{local:unit-condition}. Moreover, for every \(n\in\N_+\),
\[
 Mg(n)=h(r+D_0n),\qquad r+D_0n\ge1.
\]
\end{lemma}

\begin{proof}
Put \(u'_p=v_p(b_{m_p,p})\). These are nonnegative integers, and the
factorization argument above shows that
\(S_0=\{p:u'_p>0\}\) is finite. For \(p\in S_0\), take
\(u_p=u'_p+1\), and set
\[
 D_0=\prod_{p\in S_0}p^{u_p},
 \qquad M=\prod_{p\in S_0}p^{u_pm_p+u'_p}.
\]
Choose the representative \(-D_0<r\le0\) satisfying
\(r\equiv\mathfrak z_p\pmod{p^{u_p}}\) for every \(p\in S_0\), by the
Chinese remainder theorem. If \(S_0\) is empty, these definitions give
\(D_0=M=1\) and \(r=0\).

For \(p\in S_0\), the congruence defining \(r\) shows that
\(w_p=(\mathfrak z_p-r)/D_0\) lies in \(\Z_p\); for \(p\notin S_0\), the
same conclusion follows because \(D_0\) is a \(p\)-adic unit. Define \(g\)
initially as the displayed rational polynomial. Its centered expansion is
\[
 g(w_p+X)=\sum_{j=m_p}^k b_{j,p}\frac{D_0^j}{M}X^j.
\]
If \(p\in S_0\), the coefficient at \(j=m_p\) has valuation
\(u'_p+u_pm_p-(u_pm_p+u'_p)=0\). For \(j>m_p\), its valuation is at least
\[
 u_p(j-m_p)-u'_p\ge1.
\]
If \(p\notin S_0\), both \(D_0\) and \(M\) are \(p\)-adic units, so every
coefficient is integral and the coefficient at \(j=m_p\) is again a unit.
Thus the centered expansion belongs to \(\Z_p[X]\) for every prime \(p\).
Translation by \(-w_p\in\Z_p\) shows that \(g\) itself belongs to
\(\Z_p[x]\). Every rational coefficient of \(g\) is therefore integral at
every prime and hence is an integer.

The elements \(w_p\) are roots of \(g\); reducing them at the prime-power
factors of any modulus and applying the Chinese remainder theorem proves
intersectivity. The degree and sign of the leading coefficient are preserved
because \(D_0,M>0\). Finally, the defining identity for \(g\) gives
\(Mg(n)=h(r+D_0n)\), and \(r>-D_0\) gives \(r+D_0n\ge1\) for \(n\ge1\).
\end{proof}

In the coordinates of any progression of step \(M\), a nonzero difference
\(g(n)\), \(n\ge1\), therefore gives the nonzero original difference
\(h(r+D_0n)\) at a positive argument. This will transfer the final bound for
\(g\) back to the original polynomial. We continue below with the arbitrary
fixed polynomial \(h\) and its chosen roots, assuming
Condition~\eqref{local:unit-condition} only when explicitly stated.

Define a completely multiplicative function \(\lambda:\N_+\to\N_+\) by
\(\lambda(p)=p^{m_p}\). For each \(l\ge1\), let \(r_l\) be the unique
integer in \((-l,0]\) satisfying
\[
 r_l\equiv\mathfrak z_p\pmod{p^{v_p(l)}}
 \qquad(p\mid l),
\]
and set
\[
 h_l(x)=\frac{h(r_l+lx)}{\lambda(l)},
 \qquad a_l=\frac{a l^k}{\lambda(l)}.
\]
For \(l=1\), this convention gives \(r_1=0\) and \(h_1=h\).

\begin{lemma}[The auxiliary family]\label{local:family}
For every \(l\ge1\), the polynomial \(h_l\) belongs to \(\Z[x]\), is
intersective, and has degree \(k\) and positive leading coefficient \(a_l\).
Every coefficient of \(h_l\) is \(O_h(a_l)\). For every prime \(p\), the
element \(\mathfrak z_{p,l}=(\mathfrak z_p-r_l)/l\) belongs to \(\Z_p\), and
\begin{equation}\label{local:centered-expansion}
 \begin{aligned}
 h_l(\mathfrak z_{p,l}+X)
   &=\sum_{j=m_p}^k b_{j,p}\frac{l^j}{\lambda(l)}X^j,\\
 v_p\!\left(b_{m_p,p}\frac{l^{m_p}}{\lambda(l)}\right)
   &=v_p(b_{m_p,p}).
 \end{aligned}
\end{equation}
In particular, the root multiplicity remains \(m_p\), and
Condition~\eqref{local:unit-condition}, when it holds for \(h\), holds for
every \(h_l\) at these roots. For all positive integers \(l,D\), there is an
integer \(i\), \(0\le i<D\), such that
\begin{equation}\label{eq:1}
 \begin{gathered}
 h_l(Dx-i)=\lambda(D)h_{lD}(x),\\
 D\mid\lambda(D),\qquad \lambda(D)\le D^k,
 \qquad \rad(\lambda(D))=\rad(D).
 \end{gathered}
\end{equation}
In particular, the argument \(Dx-i\) is positive whenever \(x\in\N_+\).
\end{lemma}

\begin{proof}
Only primes dividing \(\lambda(l)\) can occur in denominators of the
coefficients of \(h_l\). Fix such a prime, put \(u=v_p(l)>0\), and use the
affine coordinate
\[
 Y=\frac{r_l-\mathfrak z_p}{p^u}+\frac{l}{p^u}x
 \quad\text{over }\Z_p.
\]
Its slope is a unit, and expansion at \(\mathfrak z_p\) gives
\begin{equation}\label{local:normalized-expansion}
 h_l(x)=
 \left(\frac{\lambda(l)}{p^{u m_p}}\right)^{-1}
 \sum_{j=m_p}^k b_{j,p}p^{u(j-m_p)}Y^j.
\end{equation}
The prefactor is a \(p\)-adic unit, and the displayed polynomial is integral
over \(\Z_p\). This proves \(h_l\in\Z[x]\).

For every prime \(p\), the element
\((\mathfrak z_p-r_l)/l\) belongs to \(\Z_p\): when \(p\mid l\), this
uses the defining congruence for \(r_l\), and otherwise \(l\) is a unit.
It is a root of \(h_l\). Reducing these roots at the prime-power factors of
an arbitrary modulus and applying the Chinese remainder theorem proves
intersectivity.

Substitution of \(\mathfrak z_{p,l}+X\) into the definition of \(h_l\)
gives the centered expansion in \eqref{local:centered-expansion}. Complete
multiplicativity gives \(v_p(\lambda(l))=m_pv_p(l)\), so
\(l^{m_p}/\lambda(l)\) is a unit in \(\Z_p\). This proves the exact valuation
there, including the preservation of the multiplicity and of the unit
condition. The assertion is local at \(p\); the displayed quotient need not
be an integer.

The degree and leading coefficient follow from the definition. Expanding
\(h(r_l+lx)\) and using \(|r_l|<l\) bounds every coefficient by
\(O_h(l^k/\lambda(l))=O_h(a_l)\).

The same root choices give \(r_{lD}\equiv r_l\pmod l\). The representatives
of this class in \((-lD,0]\) are exactly
\(r_l,r_l-l,\ldots,r_l-(D-1)l\), so \(r_{lD}=r_l-il\) for some
\(0\le i<D\). Complete multiplicativity of \(\lambda\) now gives the
polynomial identity in \eqref{eq:1}. The remaining bounds follow from
\(1\le m_p\le k\), while \(\lambda(p)=p^{m_p}\) gives the equality of prime
supports. Finally, \(Dx-i\ge D-(D-1)=1\) for \(x\ge1\).
\end{proof}

Equation~\eqref{eq:1} is the inheritance rule for avoidance. If a set
contained in one progression of step \(\lambda(D)\) avoids nonzero values
of \(h_l\) as differences, then its progression coordinates avoid nonzero
values of \(h_{lD}\). Indeed a difference \(h_{lD}(m)\), \(m\ge1\), would
give the old difference
\(\lambda(D)h_{lD}(m)=h_l(Dm-i)\), with a positive argument.

\subsection{Exponential sums uniformly over the family}

For all sufficiently large primes \(p\), the reduction of \(h_l\) is
nonconstant for every \(l\). Indeed, \(b_{m_p,p}\) is a unit outside the
finite set found above, so the first coefficient in the centered expansion
\eqref{local:centered-expansion} is a unit. If \(h_l\bmod p\) were constant,
translation by the integral element \(\mathfrak z_{p,l}\) would still give a
constant polynomial modulo \(p\), a contradiction. Under
Condition~\eqref{local:unit-condition}, this argument works at every prime:
for every \(p,l\), some nonconstant coefficient of \(h_l\) is a
\(p\)-adic unit. At any such prime with \(p>k\), the reduced degree \(k'\)
lies in \(\{1,\ldots,k\}\), so \(h_l'\bmod p\) is not the zero polynomial.
The reduced degree may be smaller than the degree over the integers.

Set
\[
 \delta=\frac{1}{100k2^k},\qquad \gamma=\frac{\delta}{10}.
\]
An argument \(x\in\F_p\) is called \emph{regular for \(h_l\)} if
\(h_l'(x)\ne0\) in \(\F_p\).

\begin{lemma}[Complete polynomial sums]\label{local:sums}
There is a threshold \(P_1>\max(k,2)\), depending only on \(h\), such that
for every prime \(p\ge P_1\) and every \(l\ge1\), all but at most \(k-1\)
arguments modulo \(p\) are regular for \(F=h_l\), and at least half are
regular. If Condition~\eqref{local:unit-condition} holds, one may choose
\(P_1=P_1(k)\), uniformly over all such polynomials and chosen roots of
degree \(k\). For every integer \(u\) with \(p\nmid u\),
\begin{equation}\label{eq:2}
 \left|\E_{x\bmod p^j}e\!\left(\frac{uF(x)}{p^j}\right)\right|
 \le p^{-8\delta j}\quad(j\ge1),
 \qquad
 \left|\E_{\substack{x\bmod p\\x\text{ regular}}}
        e\!\left(\frac{uF(x)}p\right)\right|
 \le p^{-4\delta}.
\end{equation}
The second average is normalized by the number of regular arguments.
\end{lemma}

\begin{proof}
Choose an initial threshold so that the preceding nonconstancy holds and
\(p>k\). Under Condition~\eqref{local:unit-condition}, this initial threshold
can be chosen using only \(k\). The nonzero polynomial \(F'\bmod p\) has
degree at most \(k-1\), which gives the count of regular arguments; taking
\(p\ge2(k-1)\) also makes at least half of the arguments regular.

First consider the sum modulo \(p\), and let \(k'\) be the degree of
\(F\bmod p\). If \(k'=1\), the sum is zero. If \(k'\ge2\), repeated
differencing gives
\[
 \left|\E_{x\bmod p}e(uF(x)/p)\right|^{2^{k'-1}}
 \le
 \E_{b_1,\ldots,b_{k'-1}\bmod p}
 \left|\E_{x\bmod p}
 e\!\left(\frac{u\,\Delta_{b_1}\cdots\Delta_{b_{k'-1}}F(x)}p\right)
 \right|,
\]
where \(\Delta_bF(x)=F(x+b)-F(x)\). To obtain the inequality, square the
mean, express it as the mean of multiplicative differences, and iterate the
triangle inequality and Cauchy--Schwarz. The reduction modulo \(p\) of
\(u\Delta_{b_1}\cdots\Delta_{b_{k'-1}}F\) has degree at most one, with
coefficient of \(x\) equal to \(u k'!\) times the leading coefficient
of \(F\bmod p\) times \(b_1\cdots b_{k'-1}\). Its character mean vanishes
unless one increment is zero. The right-hand side is therefore at most
\((k'-1)/p\). Since \(8\delta<2^{1-k}\le2^{1-k'}\) and
\(1\le k'\le k\), increasing the threshold by an amount depending only on
\(k\) gives the first bound in \eqref{eq:2} for \(j=1\). Removing at most
\(k-1\) arguments gives
\[
 \left|\E_{\substack{x\bmod p\\x\text{ regular}}}e(uF(x)/p)\right|
 \le\frac{p}{p-k+1}
       \left(\left|\E_{x\bmod p}e(uF(x)/p)\right|+\frac{k-1}{p}\right).
\]
The first bound and another increase depending only on \(k\) give the
second bound in \eqref{eq:2}.

For \(j\ge2\), put \(n=\lfloor j/2\rfloor\). Write
\(x=x_0+p^{j-n}y\), where \(x_0\) runs modulo \(p^{j-n}\) and \(y\)
runs modulo \(p^n\). Taylor expansion modulo \(p^j\) leaves only the
linear variation, because \(2(j-n)\ge j\). The average over \(y\) vanishes
unless \(F'(x_0)\equiv0\pmod{p^n}\). We claim that these arguments occupy
at most \(k-1\) classes modulo \(p^v\), where
\(v=\lceil n/(k-1)\rceil\).

Otherwise choose \(k\) integer representatives \(x_1,\ldots,x_k\) in
distinct such classes with \(p^n\mid F'(x_i)\). Interpolation of the
degree-at-most-\(k-1\) polynomial \(F'\) gives
\[
 F'(X)=\sum_{i=1}^k F'(x_i)
       \frac{\prod_{j\ne i}(X-x_j)}{\prod_{j\ne i}(x_i-x_j)}.
\]
Each denominator has valuation at most \((k-1)(v-1)<n\), whereas the
numerator polynomial has integral coefficients. Every coefficient of every
summand is therefore divisible by \(p\) in \(\Z_p\). This contradicts
\(F'\not\equiv0\pmod p\) and proves the claim. Since \(v\le n\le j-n\),
the density of the surviving arguments is at most
\[
 (k-1)p^{-\lceil n/(k-1)\rceil}
 \le (k-1)p^{-j/(3(k-1))}.
\]
Here \(n\ge j/3\). Since \(1/(3(k-1))>8\delta\), one fixed increase of
the prime threshold absorbs the factor \(k-1\) for every \(j\ge2\).
All choices of threshold are uniform in \(l\) and \(u\); fix \(P_1\)
large enough to satisfy them all. After the initial nonconstancy condition,
every increase used only \(k\). This proves the stated degree-only choice
under Condition~\eqref{local:unit-condition}.
\end{proof}

\subsection{Admissible residues and uniform lifting}

At large primes we have many regular arguments.
At the finitely many smaller primes, the derivative may be divisible by
\(p\) at every argument. We therefore need, uniformly in \(l\), an argument
at which \(v_p(h_l'(x))\) is bounded above.
Restrictions based on derivative nonvanishing and lifting at prime powers
also appear in Rice's treatment
\cite[Sections~2.4--2.5 and~4 of the revised version]{Rice2019}.

\begin{lemma}[A derivative bound at small primes]\label{local:small-primes}
Let
\[
 V_k=(i^j)_{0\le i,j\le k-1},\qquad
 T_{k,p}=\max_{1\le j\le k}v_p(j).
\]
For every prime \(p<P_1\), one may take
\begin{equation}\label{local:small-prime-choice}
 S_p=v_p(b_{m_p,p})+T_{k,p}+v_p(\det V_k)
\end{equation}
so that for every \(l\ge1\) some integer \(x\in\{0,\ldots,k-1\}\) satisfies
\(v_p(h_l'(x))\le S_p\). In general this choice depends only on the fixed
\(h\), its chosen roots, and \(p\). Under
Condition~\eqref{local:unit-condition}, it reduces to
\[
 S_p=T_{k,p}+v_p(\det V_k),
 \qquad \det V_k=\prod_{j=1}^{k-1}j!,
\]
and hence depends only on \(k\) and \(p\).
\end{lemma}

\begin{proof}
For a nonconstant polynomial over \(\Z_p\), consider the ideal generated
by its nonconstant coefficients. An affine substitution with unit slope
preserves this ideal: substitution gives one inclusion, since a constant
cannot contribute to a nonconstant coefficient, and the inverse affine
substitution gives the reverse inclusion. Thus it preserves the minimum
valuation of the nonconstant coefficients.

Apply this observation to the translation by \(\mathfrak z_{p,l}\).
Equation~\eqref{local:centered-expansion} displays a nonconstant coefficient
of valuation \(v_p(b_{m_p,p})\), so the minimum for \(F=h_l\) is at most
that number, uniformly in \(l\). Differentiation multiplies the coefficient
of degree \(j\) by \(j\). Hence some coefficient of \(F'\) has valuation at
most \(v_p(b_{m_p,p})+T_{k,p}\).

The matrix \(V_k\) maps the coefficient vector of a polynomial of degree at
most \(k-1\) to its values at \(0,1,\ldots,k-1\). Its determinant is
\(\prod_{0\le i<j\le k-1}(j-i)=\prod_{j=1}^{k-1}j!\), and its adjugate has
integer entries. If every value \(F'(x)\), \(0\le x<k\), had valuation
greater than \(S_p\), multiplication by \(V_k^{-1}=\operatorname{adj}(V_k)/
\det V_k\) would give valuation greater than
\(S_p-v_p(\det V_k)=v_p(b_{m_p,p})+T_{k,p}\) for every coefficient of
\(F'\). This contradicts the preceding bound and proves the lemma.
\end{proof}

Fix one such \(S_p\) for each \(p<P_1\), and define
\[
 E_p=
 \begin{cases}
  2S_p+1,&p<P_1,\\
  1,&p\ge P_1.
 \end{cases}
\]
When Condition~\eqref{local:unit-condition} holds, use the choice
\(P_1=P_1(k)\) from Lemma~\ref{local:sums} and the degree-only choice of
\(S_p\) in Lemma~\ref{local:small-primes}. Then all the exponents \(E_p\)
are determined by \(k\) and \(p\).
For \(F=h_l\), call a residue \(r\bmod p^{E_p}\) \emph{admissible for
\(l\)} if it has the following form. When \(p<P_1\), require
\(r\equiv F(x)\pmod{p^{E_p}}\) for some integer \(x\) with
\(v_p(F'(x))\le S_p\). When \(p\ge P_1\), require that \(r\) be a value
of \(F\) at a regular argument modulo \(p\). There is at least one
admissible residue at every prime. A residue equal to zero is allowed.

The next elementary lifting statement explains this choice of modulus.

\begin{lemma}[Uniform lifting on an argument class]\label{local:lifting}
Let \(F\in\Z[x]\), let \(p\) be a prime, and suppose
\(s=v_p(F'(x_0))<\infty\) for an integer \(x_0\). If \(E\) is an integer
with \(E\ge2s+1\), then
\(F\) is constant modulo \(p^E\) on the argument class
\(x_0\bmod p^{E-s}\). For every integer \(j>E\), a uniform argument in
that class modulo \(p^j\) gives a uniform value among the residues modulo
\(p^j\) congruent to \(F(x_0)\bmod p^E\).
\end{lemma}

\begin{proof}
Taylor expansion over the integers gives
\[
 F(x_0+p^{E-s}y)=F(x_0)+p^E\bigl(c_1y+pC(y)\bigr),
 \qquad p\nmid c_1,\quad C\in\Z[y].
\]
Indeed the linear coefficient is \(p^EF'(x_0)/p^s\), and every term of
degree at least two is divisible by \(p^{E+1}\), since \(E-2s\ge1\).
The polynomial \(T(y)=c_1y+pC(y)\) induces a bijection modulo every power
of \(p\): its difference at two arguments is their difference times a
\(p\)-adic unit, so the induced map on each finite residue ring is
injective and hence bijective.

For a uniform argument in the stated class modulo \(p^j\), the variable
\(y\) is uniform modulo \(p^{j-E+s}\). Its reduction modulo
\(p^{j-E}\) is still uniform, with each residue occurring \(p^s\) times.
The bijectivity of \(T\) at this latter modulus proves the asserted
distribution of \(F\).
\end{proof}

For an admissible residue of \(F=h_l\), take its defining argument \(x_0\)
and put \(s=v_p(F'(x_0))\); at a large prime any integer representative of
the regular argument has \(s=0\). The chosen exponent always satisfies
\(E_p\ge2s+1\). Lemma~\ref{local:lifting} therefore supplies an argument
class of modulus \(p^{E_p-s}\), which divides \(p^{E_p}\), on which
\(F\) has precisely the required uniform lifting property.

\subsection{Admissible cycles of a fixed length}

We finally choose one cycle length that works simultaneously at every prime
and for every member of the auxiliary family. Fix an integer \(L>2\) such
that
\[
 p^{E_p}\mid L\quad(p<P_1),
 \qquad 4\delta L>2.
\]
There are only finitely many divisibility conditions, so such a choice exists.
For example, take the least multiple of \(\prod_{p<P_1}p^{E_p}\) exceeding
\(\max\{2,(2\delta)^{-1}\}\). In general \(L\) may depend on \(h\) and its
chosen roots; under
Condition~\eqref{local:unit-condition} and the degree-only choices above,
this choice of \(L\) depends only on \(k\).

\begin{lemma}[Zero-sum admissible sequences]\label{local:cycles}
For every prime \(p\) and every positive integer \(l\), there are
\(L\) admissible residues \(r_1,\ldots,r_L\bmod p^{E_p}\) such that
\(r_1+\cdots+r_L=0\bmod p^{E_p}\).
\end{lemma}

\begin{proof}
If \(p<P_1\), repeat any admissible residue \(L\) times. The sum is zero
because \(p^{E_p}\mid L\).

If \(p\ge P_1\), let \(R\) be \(h_l(x)\) for a uniform regular
argument \(x\bmod p\), and take independent copies \(R_1,\ldots,R_L\).
Fourier inversion gives
\[
 p\,\mathbb P(R_1+\cdots+R_L=0)
 =1+\sum_{u=1}^{p-1}\bigl(\E e(uR/p)\bigr)^L.
\]
By \eqref{eq:2}, the sum following \(1\) has modulus at most
\((p-1)p^{-4\delta L}<1\). The probability is therefore positive, and
any realization with sum zero supplies the required residues.
\end{proof}

Label the \(L\) vertices of a directed cycle by starting with any initial
residue and taking successive partial sums of \(r_1,\ldots,r_L\). Their
zero sum makes the final increment close the cycle. The vertices are
distinct, but their labels need not be. The later finite-field construction
may impose a larger prime threshold \(P>P_1\) while retaining the exponents
\(E_p\) just fixed. In particular, \(E_p=1\) still holds on
\(P_1\le p<P\), where the lemma already supplies a length-\(L\) sequence.
Thus increasing that threshold does not require changing the cycle length.
